\newpage
%%%%%%%%%%%%%%%%%%%%%%%%%%%%%%%%%%%%%%%%%%%%%%%%%%%%%%%%%%%%%%%%
%%%%%%%%%%%%%%%%%%%%%%%%%%%%%%%%%%%%%%%%%%%%%%%%%%%%%%%%%%%%%%%%
%%%%%%%%%%%%%%%%%%%%%%%%%% Enunciado %%%%%%%%%%%%%%%%%%%%%%%%%%%

\begin{myblock}
\phantomsection\addcontentsline{toc}{section}{Parte A | TinyLLaMA-1.1B}
\section*{Parte A | TinyLLaMA-1.1B}

    Afinar al menos un modelo Transformer en cada parte (generación y clasificación), eligiéndolo
    de Hugging Face. 

\end{myblock}

%%%%%%%%%%%%%%%%%%%%%%%%%%%%%%%%%%%%%%%%%%%%%%%%%%%%%%%%%%%%%%%%
%%%%%%%%%%%%%%%%%%%%%%%%%%%%%%%%%%%%%%%%%%%%%%%%%%%%%%%%%%%%%%%%

%%%%%%%%%%%%%%%%%%%%%%%%%%%%%%%%%%%%%%%%%%%%%%%%%%%%%%%%%%%%%%%%
%%%%%%%%%%%%%%%%%%%%%%%%%%%%%%%%%%%%%%%%%%%%%%%%%%%%%%%%%%%%%%%%

%Los resultados de este problema se encuentran en los notebooks:

% \begin{itemize}
%     \item \textbf{MLP-MIR-Smooth.ipynb}
%     \item \textbf{exploratorio-mir-tarea01.ipynb}
% \end{itemize}

\subsection{Resultados}

\begin{lstlisting}[
    style=mystyle,
    caption={Output crudo del modelo TinyLLaMA-1.1B a nivel de palabra.},
    label={lst:output-gru-word}
]
I will take my heart with my hands and I'll lead you to believe that I am hiding something You don't need to run Heads are moving up and down, you got it It looks so much warmer there" <|startsong|> We used to play pretend, give each other different names I said I fly by the dangerous bend symbol (Wait, what? Wait, what?) And now they're outside, ready to bust (And guess what?) Lean on my pride, I'm a lion Pretending to understand how you think your world is ending (Keeps comin' around, keeps comin', keeps comin' ) I found my way, right time, wrong place But now you double as a paper maker <|endsong|> Mm, ba-da-dum, ba-da-dum-bum Dry us in the sand (Mm-mm) Da da, da da da, da-da) And when your father turns to stone <|endsong|> (Hey, hey, hey, hey, hey, hey, hey, hey) Just keep it outside, keep it outside, yeah Let me know when you've had enough Quarantine Ooh, ooh You should go your way, I should try mine Yeah\end{lstlisting}
\clearpage

\begin{lstlisting}[
    style=mystyle,
    caption={Output crudo del modelo TinyLLaMA-1.1B a nivel de palabra.},
    label={lst:output-gru-word}
]
If you ever wanna see me just let me know We get bodies every day (Ooh-ooh, ooh-ooh) Just don't believe the hype I don't have a chance to say goodbye I want to stop time Pardon my delay Don't make me waste my love Tell our dad I'm sorry I check the clock, wondering what he'll pull this time I've got two faces, Blurry's the one I'm not I've been thinkin' too much He knows that it's almost over <|endsong|> I know, I know-ow-ow And now they're outside, ready to bust (And guess what?) I could take the high road (Eh, eh, eh, eh) I prayed those lights would take me home They walked up to me and when they got close You'll be in and out, out, out, man They asked me, "Son, when I grow old Ooh, ooh So move out of our way, we're pushing sideways Yeah, yeah, yeah! Take that Know I'll keep movin', know I'll keep movin' Entertain my, entertain my faith But I know that I'm goin' low Low-key
\end{lstlisting}

\begin{lstlisting}[
    style=mystyle,
    caption={Output crudo del modelo TinyLLaMA-1.1B a nivel de palabra.},
    label={lst:output-gru-word}
]
Dema will fall like the other guys (I'll fake you out) I want to be afraid, but it seems that these days It's a backslide But they say a system's coming in In time, I will leave the city My friends and I have problems If I keep moving, they won't know You don't need to run But I can hear my way around Down there, it's some shelter from the weather They think this thing is a highway, highway I wanna see it all, no surprises, yeah Just so I don't fall through the floor So high, so high We used to play pretend, give each other different names Yeah, yeah, yeah! Don't forget abou-bou-bou-bou-bout me And now I just sit in silence To scrape the frosted windshield like we're barely scrapin' by Push on through Ah-ah-ah-ah-ah-ah, ah-ah-ah-ah-ah-ah, ah-ah-ah-ah-ah-ah (ah-ah-ah-ah-ah-ah) Can you save? Can you save my? Save my Are you doin' good? No, not me, it's for a friend (Ooh) Do-do-do
\end{lstlisting}

\subsection{Conclusiones}

El modelo TinyLLaMA-1.1B, a diferencia de las arquitecturas recurrentes tradicionales, mostró 
un desempeño significativamente superior en términos de perplejidad y calidad de los textos generados.
Tras un proceso de \textit{fine-tuning} ligero, con cuatro épocas y un \textit{batch size} de dos unidades,
la pérdida de entrenamiento descendió hasta valores cercanos a 0.61, lo cual corresponde
con una PPL de aproximadamente 1.84. En el conjunto de validación, la PPL se estabilizó 
en torno a un valor de 3.06, un valor muy bajo considerando la naturaleza del corpus. Estos resultados
superan claramente lo obtenido por las arquitecturas anteriores (RNN, LSTM y GRU), especialmente
a nivel de estabilidad y convergencia. 

En cuanto a la generación de letras, TinyLLaMA fue capaz de articular secuencias más largas y coherentes,
manteniendo tanto la estructura formal de las canciones como ciertos \textit{motifs} caracterísitcos
de Twenty One Pilots. Incrustó también conceptos que habitan el imaginario de la banda, como lo es
\textit{Blurryface}. A diferencia de las RNN, que tendieron a memorizar secuencias cortas, o las
GRU que mostraban fragmentación, TinyLLaMA logró mantener consistencia global incluso a lo largo
de estrofas enteras. 

Sin embargo, el moedlo también tiene limitaciones evidentes. Algunas secuencias derivaron en construcciones
fragmentadas o clichés genéricos como: ``\textit{I found my way, right time, wrong plac}'', una
sentencia importante y reconocida de la canción The Judge. Esto refleja que, a pesar de que 
el modelo aprovecha su capacidad de atención, sigue condicionado al tamaño relativamente pequeño
del corpus. La creatividad con la que se escribieron las canciones sigue siendo un tanto limitada, 
con muchas repeticiones que alteran la fluidez semántica por patrones rítmicos. 

En suma, podemos afirmar que TinyLLaMA-1.1B representa un salto claro en calidad frente a las arquitecturas
recurrentes. Su capacidad de auto-atención le permite modelar dependencias de largo alcance,
manteniendo tanto la coherencia semántica como los patrones rítmicos propios de las letras. Aunque
aún presenta clichés y algunas secuencias inconexas, la combinación de baja PPL, coherencia global y
estilo lírico consistente lo posicionan como un modelo superior. Quizás un corpus más extendo
y un ajuste más fino pueda superar ampliamente a esta implementación. 

Hablando en términos más coloquiales, es claro que incluso un Transformer básico y pequeño supera 
a las arquitecturas recurrentes, demostrando cómo es que este paradigma cambió la manera en que 
se hace el NLP de alto nivel. 

\clearpage